\section{Introduction}
\label{sec:introduction}

Scientific simulation, computer graphics, collision detection, and many other fields utilize data structures called octrees. An octree allows a three dimensional space to be broken down hierarchically into a tree structure with up to eight children per node. A common application would be, for example, breaking down a domain of particles, especially when paired with Barnes-Hut or Fast Multipole Methods which allow reductions in the runtime of force calculations between points from $O(n^2)$ to $O(n\log(n))$ and quasi-$O(n)$ respectively. This has driven the development of octree construction methods that operate on GPUs and across multiple nodes in clusters and supercomputers. 

One of the most interesting new methods for point clouds is Cornerstone \cite{cornerstone}. In the case of simulating the movement of particles in dynamic systems, the octree must be recalculated at each time step, however, the primary advantage of Cornerstone is to rebalance the tree rather than rebuild from scratch. Cornerstone has been heavily used on exascale supercomputers, however not all problems which may benefit from this method require such systems. Furthermore, a profile of the code internals and how the method performs across varying conditions of particle distributions, parameters, and other variants does not appear publicly but may help to optimize and evidence the method's usefulness across various systems and problem spaces. 

This project performs such a characterization and will focus closely on the mapping of Cornerstone to the current GPU architectures examining any kernel level metrics of interest, memory movement at the intra and inter node scope, as well as the interactions of communication and computation cost with different algorithm parameters.
